\chapter{Literature Review}

\section{Introduction}
The purpose of this chapter is to review the existing research and established techniques on which the proposed system is built, and to show where the present project stands in relation to them. The system combines several areas that are usually studied separately: multi-tenant software that serves many organisations from one installation, access control and authentication for sensitive records, service-based software architecture, machine learning for predicting student performance and dropout, retrieval-based question answering with large language models, and speech recognition and speech synthesis for low-resource languages such as Bangla.

Reviewing these areas serves three purposes. First, it explains the accepted methods that the project adopts, such as role-based access control, token-based sessions, password hashing, Random Forest classification, ridge regression, retrieval-augmented generation and neural speech models. Second, it identifies the limitations reported in the literature, for example the tendency of language models to produce unsupported answers and the dependence of prediction models on representative data. Third, it identifies the gap that this project addresses: these techniques are mostly presented as separate research results, while an institution needs them working together inside one secure system that is used every day. The chapter first reviews the related work by topic and then summarises how the reviewed work forms the foundation of the proposed system.

\section{Related Work}

\subsection{Multi-tenant platforms and service-based architecture}
Software as a Service (SaaS) applications commonly use multi-tenancy, in which one running application and one database serve several customer organisations while keeping their data separate. Bezemer and Zaidman \cite{bezemer2010multitenant} describe the benefits of this approach in lower hardware and maintenance cost, and they also warn that it increases the risk of one tenant's data becoming visible to another and makes maintenance more complex when tenants need different features. Database-level protection has been added to address the first risk; PostgreSQL, for example, provides row security policies that restrict which rows a database login can read or change \cite{postgres_rls}.

At the architectural level, Dragoni et al. \cite{dragoni2017microservices} review the move from single monolithic applications towards small, independently deployable services. They report better separation of concerns and independent scaling, but also note added complexity in communication, security between services and data consistency.

The proposed system follows these findings in a practical way. It uses a shared database in which every record carries an institution identifier and every query is limited to the signed-in user's institution, and it uses row security policies for the call transcript tables written by the Voice service under its own restricted database login. Instead of splitting every module into a separate service, it keeps the business modules together in one Core API and separates only the work that needs a different technology stack, namely machine learning prediction and language and speech processing. The Core API remains the only entry point for the frontend, and the other services are called with a shared service token. This keeps the security benefits of a single gateway while avoiding much of the communication overhead described in \cite{dragoni2017microservices}.

\subsection{Access control and authentication}
Role-based access control (RBAC), formalised by Sandhu et al. \cite{sandhu1996rbac}, grants permissions to roles rather than to individual users, and users receive permissions by being assigned to roles. The model simplifies administration in organisations where many people share the same duties, which matches the structure of an educational institution with administrators, teachers, accountants, wardens, storekeepers, students and guardians.

For session management, the JSON Web Token (JWT) standard \cite{rfc7519jwt} defines a compact, signed token that carries claims about the user and can be verified without a server-side session store. A known limitation of pure token-based sessions is that a token stays valid until it expires even if the user's account or role changes. For password storage, Provos and Mazi\`{e}res \cite{provos1999bcrypt} introduced bcrypt, an adaptive hashing scheme whose cost can be increased over time so that offline password guessing remains expensive.

The proposed system applies all three techniques: built-in and institution-defined roles with permission keys, signed JWT sessions, and bcrypt password hashes. It also addresses the token limitation noted above by loading the user, role and institution status from the database on every request, so a disabled account, a changed role or a suspended institution takes effect immediately. Additional measures from common security practice are included, such as locking an account after five failed sign-in attempts, short-lived tokens for users who must replace a temporary password, and an audit log of important actions.

\subsection{Machine learning for student performance and dropout prediction}
Educational data mining applies data mining and machine learning to data produced in educational settings. Romero and Ventura \cite{romero2010edm} review the field and identify prediction of student performance as one of its main tasks, together with the difficulty of moving research results into tools that institutions use in daily work. M\'{a}rquez-Vera et al. \cite{marquezvera2016dropout} show that dropout can be predicted early from academic and social indicators, and they stress the value of early warnings that allow staff to act before a student leaves.

Among the methods used for such prediction, the Random Forest classifier proposed by Breiman \cite{breiman2001randomforests} combines many decision trees trained on random subsets of data and features, which gives robust performance and provides class probabilities. For predicting a continuous value such as a grade point average, ridge regression \cite{hoerl1970ridge} adds a penalty on the size of the coefficients and gives stable estimates when the input variables are correlated. Both methods are available in the scikit-learn library \cite{pedregosa2011sklearn}.

The proposed AI service uses a Random Forest classifier over attendance rate, midterm average, assignment submission rate, previous CGPA and days of overdue fees to estimate a student's risk probability, and a ridge regression model to forecast term GPA and the resulting CGPA. Following the emphasis on actionable early warnings in \cite{marquezvera2016dropout}, the service returns the risk level together with the factors that raised it and recommended actions for staff. The reviewed studies train and evaluate their models on real student records. In contrast, the models in the current version of this project are trained on synthetic data generated when the service starts, so they demonstrate the integration of prediction into the platform rather than a validated prediction accuracy. Training on real institutional data is left as future work.

\subsection{Retrieval-based question answering with language models}
Large language models can answer questions in natural language, but they may produce fluent statements that are not supported by any source. Ji et al. \cite{ji2023hallucination} survey this problem, known as hallucination, and identify grounding answers in retrieved evidence as one of the main ways to reduce it. Lewis et al. \cite{lewis2020rag} introduced retrieval-augmented generation (RAG), in which relevant passages are first retrieved from a document collection and the model then generates an answer based on them, which also makes it possible to show the sources used.

Retrieval quality depends on the text embeddings used to compare questions with passages. Chen et al. \cite{chen2024bgem3} present BGE M3-Embedding, a multilingual embedding model that supports more than one hundred languages and long inputs, which makes it suitable for collections that mix Bangla and English. Open-weight language models such as Qwen2.5 \cite{qwen2024qwen25} can be run on an institution's own server, which keeps questions and documents inside the institution.

The proposed message assistant builds on these works but adds controls that the reviewed research does not address directly for an institutional setting. Questions about the asker's own records are answered from the ERP database rather than by the language model. Requests to change data are refused. General questions are answered by the locally hosted qwen2.5:7b model only from retrieved passages of institution documents and administrator-approved answers, which are searched by word overlap and bge-m3 embedding similarity, and the model must say when the passages do not answer the question. When no trusted answer is found, the question is not answered by guesswork but forwarded to a responsible person, and the person's answer becomes part of the knowledge base only after an administrator approves it. This human review step responds directly to the hallucination risk described in \cite{ji2023hallucination}.

\subsection{Speech recognition and speech synthesis}
Radford et al. \cite{radford2023whisper} present Whisper, a speech recognition model trained on a very large multilingual dataset with weak supervision, which recognises speech in many languages, including Bangla, and detects the spoken language. For speech synthesis, Kim et al. \cite{kim2021vits} propose VITS, an end-to-end text-to-speech model that produces natural speech directly from text. Pratap et al. \cite{pratap2023mms} extend speech technology to more than one thousand languages in the Massively Multilingual Speech (MMS) project and release VITS-based text-to-speech models for many low-resource languages, including Bangla.

The proposed Voice service uses faster-whisper, an efficient implementation of Whisper that runs on a standard processor, to convert the caller's speech to text, and the MMS Bangla and English voices to read the answer aloud. The spoken question then follows the same answering path as a typed question, so voice users receive the same verified answers and the same escalation to staff. Each call and each spoken turn is stored, which allows administrators to review the conversation later.

\subsection{Research gap}
The reviewed literature provides mature methods for each part of the problem: multi-tenancy and service-based architecture for hosting many institutions, RBAC, JWT and bcrypt for secure access, Random Forest and ridge regression for academic prediction, retrieval-augmented generation for grounded answers, and Whisper and MMS for speech in Bangla and English. However, these works are presented as separate solutions, and several of them point out the difficulty of applying research results in the daily operation of institutions \cite{romero2010edm} and the risk of unsupported answers from language models \cite{ji2023hallucination}.

The present project addresses this gap by combining these methods in one multi-institution platform that already holds the institution's academic and operational records. The assistant uses those records, the institution's own documents and administrator-approved answers, runs on a locally hosted model, supports both text and voice in Bangla and English, and passes unanswered questions to people instead of guessing. The prediction models are placed behind the same secure gateway and permission checks as the rest of the system. The contribution of the project therefore lies in the integration of established methods into a secure and usable institutional system, rather than in a new algorithm.

\section{Summary}
This chapter reviewed the research and established techniques related to the proposed system. It covered multi-tenant platforms and service-based architecture, role-based access control, token-based authentication and password hashing, machine learning for student performance and dropout prediction, retrieval-augmented question answering with language models, and speech recognition and synthesis for Bangla and English. For each area it explained how the project applies the reviewed methods and where it differs from them. The review showed that these methods are mature but usually presented separately, and that unsupported answers from language models and the gap between research models and daily institutional use remain concerns. These findings form the foundation for the proposed system, which brings the methods together in one secure platform with human review of AI answers, as described in the following chapters.
